% ==================== 7.1 工作总结 ====================
\section{工作总结}
\label{sec:summary}

本文针对月面巡视器自主行走面临的月面环境复杂多变、轮-壤相互作用机理不清、感知导航能力受限等核心难题，系统开展了面向月面巡视器自主行走的数字孪生技术研究。主要工作总结如下：

\textbf{（1）提出了面向月面巡视器自主行走的增强五维数字孪生模型}

分析了经典数字孪生五维模型在月面巡视器自主行走场景下的局限性，提出了增强的五维模型。该模型在物理实体维度融入月面环境因素，形成"巡视器-环境"一体化系统；在虚拟实体维度增加自主演化机制，提高数据时延条件下的预测能力；在服务维度增加自主决策支持能力；在孪生数据维度实现多源异构数据融合；在连接维度引入时延补偿机制。在此基础上，构建了月面巡视器自主行走数字孪生的总体架构，为后续研究提供了统一的理论框架。

\textbf{（2）建立了月面环境数字孪生模型}

针对月面地形、月壤、光照温度环境的建模需求，分别提出了相应的数字孪生建模方法。地形建模方面，基于LOLA数据构建了三维地形模型，实现了坡度、粗糙度等特征提取和障碍物识别；月壤建模方面，建立了基于Drucker-Prager准则的弹塑性本构模型和离散元模型，实现了参数分布预测；光照温度建模方面，建立了基于光线追踪的光照场模型和基于热平衡方程的温度场模型。实验验证了各模型的有效性，为巡视器动力学建模和路径规划提供了环境基础。

\textbf{（3）构建了月面巡视器动力学数字孪生模型}

研究了轮-壤交互力学理论和多体动力学建模方法，建立了考虑滑移特性和低重力效应的巡视器动力学模型。通过地面模拟实验验证了模型的有效性，仿真结果与实验数据的误差在10\%以内。构建了动力学数字孪生仿真系统，实现了实时仿真，为路径规划和避障决策提供了动力学支撑。

\textbf{（4）提出了基于数字孪生的月面岩石识别与避障方法}

构建了数字孪生平台下的岩石识别框架，开发了增强型ORB-SLAM2系统，实现了月面岩石的精确识别与定位。设计了基于威胁评估的避障策略，包括停止、绕行、越过等多种避障方式。实验结果表明，岩石检测mAP达到94.0\%，定位误差90\%小于0.5 m，避障成功率超过96\%。

\textbf{（5）研究了基于数字孪生的巡视器路径规划方法}

构建了数字孪生环境下的分层路径规划框架，提出了基于A-D3QN的混合路径规划算法。该算法结合Double DQN、Dueling网络结构和注意力机制，实现了全局规划与局部规划的协同优化。实验结果表明，在复杂场景中路径长度缩短、成功率提升、能耗降低，满足月面巡视器自主行走的实际需求。
